% !TEX root = main.tex
\section{Material and methods}

\subsection{System overview}
The chamber is an acrylic box of approximately 100\,×\,100\,×\,80\,mm with three 12\,V/5\,W W5W
lamps on the lid and an aluminum plate as thermal mass. An ESP32 DevKit v1 board is the only
processing element: it reads the temperature and the battery voltage, computes the control action
and drives the lamps through an IRLZ44N logic-level MOSFET~\cite{irlz44n}, as shown in
Figure~\ref{fig:mcu}. The system is powered by a 3S 18650 battery pack, with a buck converter
supplying 5\,V to the board.

\begin{figure}[htb]
  \centering
  \resizebox{\linewidth}{!}{%
  \begin{tikzpicture}[font=\footnotesize,
      b/.style={draw, rounded corners=2pt, fill=gray!12, align=center,
                minimum height=0.9cm, minimum width=1.9cm},
      seta/.style={-{Latex[length=1.6mm]}, thick}]
    \node[draw, thick, rounded corners=3pt, fill=blue!6, align=center,
          minimum width=3.2cm, minimum height=3.3cm] (mcu)
          {\textbf{ESP32}\\[2pt]ADC1 (12 bits)\\LEDC (PWM)\\Timer / sleep\\Wi-Fi};
    \node[b, left=0.9cm of mcu.north west, anchor=north east] (lm) {LM35\\10\,mV/°C};
    \node[b, left=0.9cm of mcu.south west, anchor=south east] (vb) {Battery\\divider};
    \node[b, right=0.9cm of mcu.north east, anchor=north west] (mos) {IRLZ44N +\\3$\times$W5W};
    \node[b, right=0.9cm of mcu.south east, anchor=south west] (web) {Web\\interface};
    \draw[seta] (lm.east) -- (lm.east -| mcu.west);
    \draw[seta] (vb.east) -- (vb.east -| mcu.west);
    \draw[seta] (mos.west -| mcu.east) -- node[above]{PWM} (mos.west);
    \draw[seta, <->] (web.west -| mcu.east) -- node[above]{Wi-Fi} (web.west);
  \end{tikzpicture}}
  \caption{Microcontroller-centered architecture of the thermal chamber.}
  \label{fig:mcu}
\end{figure}

\subsection{Analog acquisition}
The LM35 sensor outputs 10\,mV/°C~\cite{lm35} and is connected to an ADC1 channel. ADC1 was chosen
because, on the ESP32, ADC2 cannot be used while the Wi-Fi driver is active~\cite{esp32trm}. Each
reading is the average of 16 consecutive conversions, which reduces noise without a hardware
filter. The battery voltage, needed for the power compensation, is read through a resistive divider
on a second ADC1 channel \pendente{inform the divider resistors and the ADC attenuation used}.

\subsection{PWM generation}
The lamps are driven by the LEDC (LED control) peripheral configured at 500\,Hz with 10-bit
resolution (1024 duty levels). The low frequency keeps the MOSFET switching losses small, and the
thermal inertia of the chamber makes the PWM ripple irrelevant to the temperature. The LEDC timer
must be clocked by a source that remains active in light sleep, so that the lamps keep receiving the
PWM signal while the processor sleeps \pendente{confirm the LEDC clock source used in the firmware}.

\subsection{Control algorithm}
At each control period the firmware computes the error between the setpoint and the measured
temperature and updates a PI controller~\cite{astrom1995}:
\begin{equation}
  u = K_p\,e + K_i \int e\,dt
  \label{eq:pi}
\end{equation}
The integral term is frozen while the output is saturated (anti-windup). Because the lamp power is
proportional to the square of the supply voltage, the duty cycle is corrected by the measured
battery voltage:
\begin{equation}
  D_{corr} = D_{PI}\left(\frac{12}{V_{bat}}\right)^{2}
  \label{eq:comp}
\end{equation}

\subsection{Operating modes and power management}
The firmware has two modes. In \emph{bench mode}, Wi-Fi is on, the loop runs every 1\,s and a web
page shows the setpoint, temperature and duty cycle in real time. In \emph{embedded mode}, Wi-Fi is
off and the ESP32 enters light sleep between control updates, taken every 60\,s. Light sleep was
chosen instead of deep sleep because deep sleep turns off the digital peripherals and would stop
the PWM output, whereas light sleep preserves the peripheral state and resumes execution where it
stopped~\cite{esp32trm}.

\begin{figure}[htb]
  \centering
  \begin{tikzpicture}[font=\footnotesize, node distance=0.25cm,
      p/.style={draw, rounded corners=2pt, fill=blue!6, align=center, minimum width=4.6cm, minimum height=0.5cm},
      seta/.style={-{Latex[length=1.5mm]}, thick}]
    \node[p] (w) {Wake up (timer)};
    \node[p, below=of w] (a) {Read LM35 and $V_{bat}$ (16 samples)};
    \node[p, below=of a] (c) {PI + anti-windup};
    \node[p, below=of c] (d) {Battery compensation};
    \node[p, below=of d] (e) {Update LEDC duty};
    \node[p, below=of e] (s) {Light sleep (60\,s)};
    \foreach \x/\y in {w/a, a/c, c/d, d/e, e/s} \draw[seta] (\x) -- (\y);
    \draw[seta] (s.east) -- ++(0.5,0) |- (w.east);
  \end{tikzpicture}
  \caption{Main loop of the firmware in embedded mode.}
  \label{fig:loop}
\end{figure}

\subsection{Experimental procedure}
\pendente{Describe the tests: current consumption of the board in bench mode, in embedded mode
awake and in light sleep (multimeter or shunt); estimated and measured battery runtime; and the
closed-loop response to a setpoint in each mode.}
